\documentclass[11pt]{article}

\usepackage[T1]{fontenc}
\usepackage[margin=1in]{geometry}
\usepackage{amsmath,amssymb,amsthm}
\usepackage{aliascnt}
\usepackage{booktabs}
\usepackage{graphicx}
\usepackage{microtype}
\usepackage[hidelinks]{hyperref}
\usepackage[capitalise,noabbrev]{cleveref}

\newcommand{\C}{\mathcal G}
\newcommand{\E}{\mathbb E}
\newcommand{\Pp}{\mathbb P}
\newcommand{\Aut}{\operatorname{Aut}}
\newcommand{\Spec}{\operatorname{Spec}}
\newcommand{\Pois}{\operatorname{Pois}}
\newcommand{\TV}{\mathrm{TV}}

\theoremstyle{plain}
\newtheorem{theorem}{Theorem}[section]
\newaliascnt{proposition}{theorem}
\newtheorem{proposition}[proposition]{Proposition}
\aliascntresetthe{proposition}
\newaliascnt{lemma}{theorem}
\newtheorem{lemma}[lemma]{Lemma}
\aliascntresetthe{lemma}
\newaliascnt{corollary}{theorem}
\newtheorem{corollary}[corollary]{Corollary}
\aliascntresetthe{corollary}
\theoremstyle{definition}
\newaliascnt{definition}{theorem}
\newtheorem{definition}[definition]{Definition}
\aliascntresetthe{definition}
\newaliascnt{algorithm}{theorem}
\newtheorem{algorithm}[algorithm]{Algorithm}
\aliascntresetthe{algorithm}
\newaliascnt{remark}{theorem}
\newtheorem{remark}[remark]{Remark}
\aliascntresetthe{remark}

\numberwithin{equation}{section}

\title{The Local Whitehead-Move Chain on Cubic Multigraphs}
\author{Alexander Kolpakov\\[0.5ex]
\small Center for Science, Technology, Engineering, and Mathematics\\
\small University of Austin, Austin, Texas, United States\\
\small \texttt{akolpakov@uaustin.org}\\[0.75ex]
\small The Wolfram Institute for Computational Foundations of Science\\
\small \texttt{akolpakov@wolframinstitute.org}}
\date{}

\begin{document}
\maketitle

\begin{abstract}
We study the random walk on the graph of connected cubic multigraphs under
Whitehead moves.  For even $N\geq2$, the states are isomorphism classes of
graphs with $N$ vertices.  At each step, an edge and one of the two Whitehead
moves on it are chosen uniformly.  Loops and multiple edges are allowed, and
both moves on a loop act trivially.  Moves with isomorphic results are counted
with multiplicity.

We prove that the chain is reversible and determine its stationary
distribution: the cubic configuration model conditioned on connectedness,
viewed on isomorphism classes.  For its transition matrix $P_N$, we show that
$\Spec(P_N)\subseteq[-1/2,1]$.  Functions of the number of loops give, for
every fixed $r$, the bound $2r/(3N)+o(N^{-1})$ on the $(r+1)$st smallest
eigenvalue of $I-P_N$.  We also prove that the cut between bridged and
bridgeless graphs has conductance exactly $1/(3N/2-1)$.  These estimates give
linear lower bounds for relaxation and mixing times.

An exhaustive enumeration gives the full transition matrices for every even
$N$ from $2$ through $14$, reaching $21{,}096$ classes and $886{,}032$
Whitehead moves at $N=14$.  The computations verify the stationary weights,
bridge flow, and exact conductance formula.
\end{abstract}

\noindent\textbf{Keywords:} Whitehead move, cubic multigraph, configuration
model, Markov chain, conductance, reconfiguration graph.

\noindent\textbf{2020 Mathematics Subject Classification:}
05C81, 05C40, 60J10, 05C30.

\section{Introduction}

The graph of cubic graphs has isomorphism classes of connected cubic
multigraphs as vertices and Whitehead moves as directed edges.  Such a move
collapses a non-loop edge and expands the resulting four-valent vertex in one
of two other ways.  We allow loops and multiple edges, and a move on a loop is
defined to leave the graph unchanged.  There is one directed edge for each
move, even when several moves have isomorphic results.

Whitehead moves also arise from pants decompositions.  When $N=2g-2$, a cubic
multigraph can be used as the dual graph of a pants decomposition of a closed
surface of genus $g$: put one vertex in each pair of pants and one edge across
each glued cuff.  The surface retains gluing data that the abstract graph
forgets.  Hatcher describes the corresponding complex of pants decompositions
\cite{Hatcher1999}; Rafi and Tao relate Whitehead moves on trivalent graphs to
the geometry of moduli space \cite{RafiTao2013}.

Grave de Peralta and Kolpakov proved that this directed graph is connected
and that its outer conductance tends to zero as the number of vertices tends
to infinity \cite{GraveDePeraltaKolpakov2024}.  Their conductance is defined
using outgoing-edge volume.  Each vertex has $3N$ outgoing edges, counted
with multiplicity, so this volume is proportional to the number of classes
in a set.  Li studies the Laplacians of the simple labelled and unlabelled
graphs obtained by identifying repeated edges \cite{Li2024}.

Here we consider the random walk that chooses an edge of the cubic graph and
one of its two Whitehead moves uniformly.  Its transition matrix is obtained
by dividing the directed adjacency matrix by $3N$.  Although each move is
reversible, the numbers of moves between two isomorphism classes need not
agree in the two directions.  In particular, the uniform distribution on
classes need not be stationary.  We determine the stationary distribution
and use it to study conductance, the spectrum, and convergence to equilibrium.

The stationary distribution is obtained by lifting the walk to labelled
half-edge configurations, where each move is a symmetric exchange.  The
spectral estimates follow from two local calculations.  Fixing the chosen
edge gives a walk on the Johnson graph $J(4,2)$, while counting the moves that
remove a loop gives a Dirichlet-form identity for functions of the number of
loops.  For the conductance estimate, we show that the only moves from a
bridged graph to a bridgeless graph are the two moves on its unique bridge.

More precisely, let $m_{ii}$ count the loops at vertex $i$, let $m_{ij}$ count
the edges between distinct vertices $i$ and $j$, and let $\Aut_v(G)$ be the
group of vertex permutations preserving these multiplicities.  The stationary
weight of the class of $G$ is proportional to
\begin{equation}\label{eq:intro-weight}
  \frac{1}{|\Aut_v(G)|}
  \frac{(3!)^N}
  {\prod_i 2^{m_{ii}}m_{ii}!\prod_{i<j}m_{ij}!}.
\end{equation}
We prove that $\Spec(P_N)\subseteq[-1/2,1]$ and that, for every fixed $r$,
\[
  \lambda_{r+1}(I-P_N)\leq \frac{2r}{3N}+o(N^{-1}).
\]
For the bridged--bridgeless cut we obtain conductance
$1/(3N/2-1)$ for every even $N\geq2$.  The proof uses the formal
generating functions of connected and bridgeless cubic graphs, related by
the combinatorial Legendre transform \cite[Section~5.8]{Borinsky2018}.
Hence, for fixed $\varepsilon<1-e^{-1}$,
\[
  t_{\mathrm{mix}}(\varepsilon)
  \geq \frac{3N}{2}\bigl(1-e^{-1}-\varepsilon\bigr)(1+o(1)).
\]
The stationary distribution is given in \cref{thm:stationary}.  The spectral
and loop-count estimates are proved in
\cref{prop:spectral-interval,thm:loop-form}, and the conductance and mixing
bounds in \cref{thm:bridge-flow,thm:exact-conductance,cor:mixing}.  Finally,
\cref{sec:enumeration} gives an exhaustive enumeration for
$N=2,4,6,8,10,12,14$.  All limits are taken through even $N$.

\section{The Local Whitehead Chain}\label{sec:chain}

Let $\C_N$ be the set of isomorphism classes of connected cubic multigraphs
on $N$ vertices.  A loop contributes two to the degree, and parallel edges
are counted separately.

Let $e=uv$ be a non-loop edge.  Besides $e$, two half-edges $a,b$ are
incident with $u$ and two half-edges $c,d$ are incident with $v$.  The current
local partition is
\[
  \{a,b\}\mid\{c,d\}.
\]
Collapsing $e$ and expanding again gives the two other partitions
\[
  \{a,c\}\mid\{b,d\},
  \qquad
  \{a,d\}\mid\{b,c\}.
\]
These define the two Whitehead moves $w_e^{(1)}$ and $w_e^{(2)}$ on $e$;
see \cref{fig:whitehead-move}.  The symbols $a,b,c,d$ denote half-edges;
two of them may belong to the same edge when loops or multiple edges are
present.

\begin{figure}[ht]
  \centering
  \includegraphics[width=\textwidth]{fig-whitehead-move.pdf}
  \caption{The local move at a non-loop edge.  The central edge is retained,
  while the four other incident half-edges are repartitioned between its two
  endpoints.  There are exactly three partitions of a four-element set into
  two unordered pairs: the current partition and the two resolutions.}
  \label{fig:whitehead-move}
\end{figure}

If $e$ is a loop, we define both $w_e^{(1)}$ and $w_e^{(2)}$ to be the
identity, as in \cite{GraveDePeraltaKolpakov2024}.  Thus every graph has
$3N$ Whitehead moves, indexed by an edge and one of the two choices on it.
These moves are counted separately even if they produce isomorphic graphs
or leave the isomorphism class unchanged.  For $G,H\in\C_N$, let $m_{GH}$
be the number of moves taking a representative of $G$ to a graph isomorphic
to $H$.  This number does not depend on the representative, and
\[
  \sum_{H\in\C_N}m_{GH}=3N.
\]

\begin{definition}[Local Whitehead chain]
The transition matrix $P_N$ on $\C_N$ is
\begin{equation}\label{eq:kernel}
  P_N(G,H)=\frac{m_{GH}}{3N}.
\end{equation}
\end{definition}

\begin{proposition}[Connected support]\label{prop:connected-support}
The support of $P_N$ is connected for every even $N\geq2$.
\end{proposition}

\begin{proof}
This is the connectivity assertion in
\cite[Theorem~2.2]{GraveDePeraltaKolpakov2024}, proved there with the same move
and loop conventions.  We include the reduction because connectivity is also
what makes the enumeration in \cref{sec:enumeration} exhaustive.  For $N=2$,
the two classes are adjacent, as the two-state chain below shows.

Suppose $N\geq4$.  If the non-loop edges contain a cycle, choose a shortest
one, of length $k\geq2$.  It has no chord.  At one of its edges, one resolution
places the two neighboring cycle half-edges at the same endpoint of the
central edge.  The cycle is thereby shortened from length $k$ to length
$k-1$; when $k=2$, the remaining edge becomes a loop.  Continuing along the
shortened cycle turns it into a loop after $k-1$ moves.  Each completed
reduction creates one loop and removes none.  Since a cubic graph has at most
one loop at each vertex, repeated reductions eventually leave no non-loop
cycle.  The non-loop edges are then acyclic and connected, since loops do not
join different vertices.  They therefore form a tree, and every leaf of that
tree carries a loop.  There are $3N/2-(N-1)=N/2+1$ loops, so all resulting
rooted trees have the same number of leaves.

It remains to connect all such looped trees.  Choose a leaf $v$, delete all
loops, delete $v$ and its remaining incident edge, and mark the newly exposed
vertex as the root.  The result is a rooted full binary tree: after orienting
its edges away from the root, every non-leaf has two children.  Resolving an
internal edge performs an ordinary tree rotation.  Give the two children at
each vertex an arbitrary left--right order.  If the left child of the root is
not a leaf, rotate its edge with the root.  The new left subtree is a proper
subtree of the old one, so repetition eventually makes the left child a leaf.
Now leave that leaf in place and apply the same procedure inside the right
subtree.  Induction on the number of internal vertices produces the right
comb, in which every internal vertex has a leaf as its left child except at
the final cherry.  After the left--right order is forgotten, this comb depends
only on its number of leaves.  Restoring $v$ and all the loops therefore gives
a fixed looped cubic tree determined only by $N$.  Hence every class can be
joined to that class.
\end{proof}

Thus $P_N$ is irreducible.  Since $\C_N$ contains graphs with loops and the
moves on a loop leave the state unchanged, it is also aperiodic.  It therefore
has a unique stationary distribution.  In general $m_{GH}\ne m_{HG}$, so this
distribution need not be uniform.

The two states at $N=2$ already show why multiplicity matters.  Write $G_0$
for the triple edge and $G_1$ for the graph consisting of one joining edge and
one loop at each endpoint; see \cref{fig:two-state}.  Directly counting the
six moves from each graph gives
\[
  (m_{G_iG_j})=
  \begin{pmatrix}3&3\\2&4\end{pmatrix},
  \qquad
  P_2=
  \begin{pmatrix}1/2&1/2\\1/3&2/3\end{pmatrix}.
\]
The stationary probabilities are $2/5$ and $3/5$, not $1/2$ and $1/2$.

\begin{figure}[ht]
  \centering
  \includegraphics[width=.88\textwidth]{fig-two-state-chain.pdf}
  \caption{The chain for $N=2$.  Arrow labels give the numbers of Whitehead
  moves between the two classes, including moves that leave the class
  unchanged.  The transition multiplicities in the two directions differ.}
  \label{fig:two-state}
\end{figure}

\section{Configuration Space and Stationary Measure}\label{sec:stationary}

To determine the stationary distribution, give each
vertex $i\in[N]$ three distinguishable
half-edges $(i,1),(i,2),(i,3)$.  A \emph{configuration} is a perfect matching
of these $3N$ half-edges.  Contracting each triple to one vertex turns a
configuration into a labelled cubic multigraph.  A matched pair within one
triple becomes a loop.

For a labelled cubic multigraph $G$ on $[N]$, let $m_{ii}$ be the number
of loops at $i$ and $m_{ij}$ the number of edges joining $i$ and $j$ for
$i<j$.  Define
\begin{equation}\label{eq:configuration-count}
  c(G)=
  \frac{(3!)^N}
  {\prod_i2^{m_{ii}}m_{ii}!\prod_{i<j}m_{ij}!}.
\end{equation}
Let $\Aut_v(G)$ be the group of vertex permutations preserving all adjacency
multiplicities.

\begin{lemma}[Counting configurations]\label{lem:configuration-fibers}
A fixed labelled cubic multigraph $G$ is represented by exactly $c(G)$
half-edge configurations.  Consequently, the number of configurations whose
unlabelled graph lies in the class $G\in\C_N$ is
\begin{equation}\label{eq:class-fiber}
  \frac{N!}{|\Aut_v(G)|}\,c(G).
\end{equation}
\end{lemma}

\begin{proof}
At each vertex there are $3!$ ways to assign its three labelled half-edges to
the three incidences of the adjacency multigraph, giving $(3!)^N$ assignments
in all.  Each matching is counted twice for every loop, $m_{ii}!$ times for
the ordering of the loops at vertex $i$, and $m_{ij}!$ times for the ordering
of the parallel edges between $i$ and $j$.  Dividing by these factors gives
\eqref{eq:configuration-count}.

There are $N!/|\Aut_v(G)|$ labelled adjacency multigraphs in the unlabelled
class of $G$.  Multiplication by $c(G)$ proves
\eqref{eq:class-fiber}.
\end{proof}

We next define the move on configurations.  Regard each matched pair as an
edge whose ends occupy labelled positions at their incident vertices.  If the
chosen edge joins distinct vertices $u$ and $v$, choose one of the other two
edge ends at $u$ and one at $v$, and exchange their positions.  There are four
choices.  In the notation of
\cref{fig:whitehead-move}, exchanging $a$ with $c$ or $b$ with $d$ gives one
resolution, while exchanging $a$ with $d$ or $b$ with $c$ gives the other.
If the chosen edge is a loop, all four exchanges are defined to be the
identity.  There are therefore $(3N/2)\cdot4=6N$ exchanges from each
configuration, counted with multiplicity.

\begin{lemma}[Symmetric lift]\label{lem:symmetric-lift}
The lifted configuration kernel is symmetric.  Its restriction to connected
configurations is closed, and after projection to connected isomorphism
classes it is the kernel $P_N$ of \eqref{eq:kernel}.
\end{lemma}

\begin{proof}
Each exchange is an involution: its reverse exchanges the same two edge ends.
Hence the number of lifted exchanges from a configuration $x$ to a configuration
$y$ equals the number from $y$ to $x$.  Uniform measure on configurations is
therefore reversible.  Relabelling vertices or permuting the three positions at a
vertex commutes with the lifted transition, so its projection is a Markov
kernel on isomorphism classes.

Connectedness is preserved: collapse the chosen edge to obtain a connected
quotient, then expand the four-valent vertex into two vertices still joined by
the central edge.  The connected configurations therefore form a closed
subset.

For a non-loop edge, each of the two Whitehead moves has two realizations
among the four exchanges.  Its projected probability is therefore
$2/(6N)=1/(3N)$.  At a loop, all four exchanges leave the graph unchanged,
giving the same total probability as the two trivial Whitehead moves.
Summing over moves with isomorphic results gives $m_{GH}/(3N)$.
\end{proof}

\begin{theorem}[Stationary law]\label{thm:stationary}
For all $G,H\in\C_N$,
\begin{equation}\label{eq:detailed-balance-counts}
  \frac{c(G)}{|\Aut_v(G)|}m_{GH}
  =
  \frac{c(H)}{|\Aut_v(H)|}m_{HG}.
\end{equation}
Consequently, the local Whitehead chain is reversible with stationary law
\begin{equation}\label{eq:stationary}
  \pi_N(G)=
  \frac{c(G)/|\Aut_v(G)|}
  {\displaystyle\sum_{K\in\C_N}c(K)/|\Aut_v(K)|}.
\end{equation}
This is the cubic configuration model conditioned on connectedness and
projected to isomorphism classes.
\end{theorem}

\begin{proof}
By \cref{lem:symmetric-lift}, uniform measure conditioned on connected
configurations is reversible and projects to $P_N$.  The mass of an unlabelled
fibre is given by \eqref{eq:class-fiber}.  The common factor $N!$ cancels on
normalization, which gives \eqref{eq:stationary}.  Summing detailed balance
on configurations over the fibres of $G$ and $H$ gives
\eqref{eq:detailed-balance-counts}.
\end{proof}

\begin{remark}
Automorphism order alone does not determine the stationary weight.  In the
$N=2$ example, both graphs have vertex-automorphism group of order two, but
$c(G_0)=6$ and $c(G_1)=9$, producing the stationary ratio $2:3$.
\end{remark}

\section{Pairing-Model Estimates}\label{sec:pairing}

By \cref{thm:stationary}, we may use the pairing model to estimate stationary
probabilities on $\C_N$.  This model was introduced by Bollob\'as
\cite{Bollobas1980}; see also the survey \cite{Wormald1999}.  We shall use the
following estimates.

\begin{lemma}[Configuration estimates]\label{lem:configuration-estimates}
Let $\widehat G_N$ be the cubic multigraph obtained from a uniform pairing of
$3N$ half-edges in $N$ labelled boxes of size three, without conditioning on
connectedness.

\begin{enumerate}
\item $\Pp(\widehat G_N\text{ is disconnected})=O(N^{-1})$.
\item Conditional on connectedness, the number $L$ of loops converges in
distribution and in every fixed moment to $\Pois(1)$.
\item Conditional on connectedness, the probability of a bridge whose smaller
side has at least three vertices is $O(N^{-1})$.
\end{enumerate}
\end{lemma}

\begin{proof}
\noindent\emph{Disconnected configurations.}
Write $\mathcal C_N$ for the event that $\widehat G_N$ is connected.  We first
bound $\Pp(\mathcal C_N^c)$.  If a set of $k\leq N/2$ vertices is a union of
connected components, then all $3k$ half-edges in the set are paired
internally; in particular, $k$ is even.  The expected number of such
$k$-sets is
\begin{equation}\label{eq:disconnected-subsets}
  D_{N,k}=\binom Nk
  \frac{(3k-1)!!(3N-3k-1)!!}{(3N-1)!!}.
\end{equation}
Here and below $(-1)!!=1$.  The term $k=2$ can be read off directly:
\[
  D_{N,2}
  =\binom N2\frac{5!!}{(3N-1)(3N-3)(3N-5)}
  =\frac{5}{18N}+O(N^{-2}).
\]
We give the uniform estimate used to sum the remaining terms.  There are
absolute constants $a,b>0$ such that, for every $m\geq1$,
\begin{equation}\label{eq:stirling-bounds}
 a\sqrt m\,(m/e)^m\leq m!\leq b\sqrt m\,(m/e)^m.
\end{equation}
For even $M\geq2$, the identity
$(M-1)!!=M!/(2^{M/2}(M/2)!)$ therefore gives
\begin{equation}\label{eq:double-factorial-bounds}
 a_0(M/e)^{M/2}\leq (M-1)!!
 \leq b_0(M/e)^{M/2}
\end{equation}
with absolute positive constants $a_0,b_0$.  The same bounds give, uniformly
for $1\leq k<N$,
\[
 \binom Nk\leq
 C_0\sqrt{\frac{N}{k(N-k)}}\exp\{NH(k/N)\},
 \qquad
 H(x)=-x\log x-(1-x)\log(1-x).
\]
Applying these inequalities to \eqref{eq:disconnected-subsets} and collecting
the powers of $k$, $N-k$, and $N$ gives, for even $4\leq k\leq N/2$,
\begin{equation}\label{eq:disconnected-entropy-bound}
  D_{N,k}\leq
  C\sqrt{\frac{N}{k(N-k)}}
  \exp\left\{-\frac N2 H(k/N)\right\}.
\end{equation}
Indeed, writing $x=k/N$, the powers from the double factorials satisfy
\[
 \frac{(3k)^{3k/2}(3N-3k)^{(3N-3k)/2}}
      {(3N)^{3N/2}}
 =x^{3Nx/2}(1-x)^{3N(1-x)/2}
 =\exp\{-3NH(x)/2\}.
\]
The binomial coefficient contributes $\exp\{NH(x)\}$, leaving the
exponential and square-root prefactor in
\eqref{eq:disconnected-entropy-bound}.

When $k\leq N/4$, the inequality
$H(k/N)\geq(k/N)\log(N/k)$ and $N-k\geq N/2$ reduce the right-hand side to
\[
 u_k=Ck^{-1/2}(k/N)^{k/2}.
\]
For even $k\geq4$ with $k+2\leq N/4$,
\[
 \frac{u_{k+2}}{u_k}
 =\sqrt{\frac{k}{k+2}}\frac{k+2}{N}
   \left(1+\frac2k\right)^{k/2}
 \leq \frac e4<1.
\]
Thus the small-$k$ terms form a geometric tail beginning with
$u_4=O(N^{-2})$.  In the range $N/4\leq k\leq N/2$, the entropy is at least
$H(1/4)$; even after summing $O(N)$ terms, the contribution is exponentially
small.  A union bound now proves
$\Pp(\mathcal C_N^c)=O(N^{-1})$.

\smallskip
\noindent\emph{Loops.}
A vertex supports at most one loop in a cubic
configuration.  To specify an ordered $r$-tuple of loops, choose $r$ distinct
vertices, choose one of the three pairs of half-edges at each vertex, and
require those $r$ pairs to occur in the random matching.  Therefore, for
fixed $r$,
\begin{equation}\label{eq:loop-factorial-moments}
  \E (L)_r
  =(N)_r3^r
  \prod_{j=0}^{r-1}\frac{1}{3N-(2j+1)}
  \longrightarrow 1.
\end{equation}
These are the limiting factorial moments of $\Pois(1)$, a moment-determinate
law.  Since every power $L^q$ is a fixed linear combination of
$(L)_0,\ldots,(L)_q$, \eqref{eq:loop-factorial-moments} also gives
$\sup_N\E L^q<\infty$ for every fixed $q$.  To justify conditioning, fix $q$
and use Cauchy--Schwarz:
\[
  \E\bigl[L^q\mathbf 1_{\mathcal C_N^c}\bigr]
  \leq \bigl(\E L^{2q}\bigr)^{1/2}
       \Pp(\mathcal C_N^c)^{1/2}=o(1).
\]
It follows that
\[
 \left|\E[L^q\mid\mathcal C_N]-\E L^q\right|=o(1).
\]
Moreover, for every event $E$ determined by $L$,
\[
 \left|\Pp(E\mid\mathcal C_N)-\Pp(E)\right|
 \leq \frac{2\Pp(\mathcal C_N^c)}{\Pp(\mathcal C_N)}=o(1).
\]
Conditioning therefore preserves both the distributional limit and every
fixed moment.

\smallskip
\noindent\emph{Large bridge sides.}
Suppose a bridge cuts off a vertex set $S$ of size $k\leq N/2$.
The degree sum on $S$ gives $3k=2|E(S)|+1$, so $k$ is odd.  To count candidate
sets, choose one crossing half-edge on each side and pair all remaining
half-edges internally.  For odd $k\geq3$, the expected number of $k$-sets
with exactly one crossing pair is therefore
\begin{equation}\label{eq:one-crossing-subsets}
  A_{N,k}=\binom Nk(3k)(3N-3k)
  \frac{(3k-2)!!(3N-3k-2)!!}{(3N-1)!!}.
\end{equation}
The first term is
\[
  A_{N,3}=\frac{35}{18N}+O(N^{-2}).
\]
Apply \eqref{eq:stirling-bounds}--\eqref{eq:double-factorial-bounds} to
\eqref{eq:one-crossing-subsets}.  Replacing $3k-1$ by $3k$, and
$3(N-k)-1$ by $3(N-k)$, changes the result by bounded factors.  After
cancellation, the upper bound is
\[
 C\sqrt{\frac{N}{k(N-k)}}e^{NH(k/N)}k(N-k)
 \frac{(3k)^{(3k-1)/2}(3N-3k)^{(3N-3k-1)/2}}
      {(3N)^{3N/2}}.
\]
The powers of $k$ and $N-k$ outside the exponential reduce its prefactor to
$C\sqrt N$, while the entropy exponent is again $N H-3NH/2$.  Thus
\begin{equation}\label{eq:one-crossing-entropy-bound}
 A_{N,k}\leq C\sqrt N
 \exp\left\{-\frac N2H(k/N)\right\},
 \qquad 3\leq k\leq N/2.
\end{equation}
For $k\leq N/4$, this is at most
$v_k=C\sqrt N(k/N)^{k/2}$.  If $k+2\leq N/4$, then
\[
 \frac{v_{k+2}}{v_k}
 =\frac{k+2}{N}\left(1+\frac2k\right)^{k/2}
 \leq\frac e4.
\]
The odd tail beginning at
$k=5$ is therefore $O(v_5)=O(N^{-2})$.  In the central range, summing
\eqref{eq:one-crossing-entropy-bound} gives
$O(N^{3/2}e^{-NH(1/4)/2})$.  Together with the displayed value of $A_{N,3}$,
this proves that the sum is $O(N^{-1})$.
The sum of the $A_{N,k}$ is an upper bound for the probability in question,
because the internally paired sides counted in
\eqref{eq:one-crossing-subsets} need not themselves be connected.
Conditioning on $\mathcal C_N$, whose probability is $1-O(N^{-1})$, changes
the bound only by a factor $1+O(N^{-1})$.
\end{proof}

By \cref{thm:stationary}, these conditional laws are precisely the laws of the
corresponding statistics under $\pi_N$.

\section{The Local Spectrum}\label{sec:spectrum}

Since the chain is reversible, $P_N$ is self-adjoint on $L^2(\pi_N)$.
We obtain a lower bound for its spectrum by considering the four exchanges
at a fixed edge.

\begin{proposition}[Spectral interval]\label{prop:spectral-interval}
For every even $N$,
\[
  \Spec(P_N)\subseteq[-1/2,1].
\]
\end{proposition}

\begin{proof}
Let $\widetilde\Omega_N$ be the following labelled cover of the connected
configuration space.  Fix $3N/2$ labelled abstract edges, give each edge two
distinguishable ends, and assign the $3N$ ends bijectively to three
labelled positions at each of $N$ labelled vertices, retaining the assignments
whose induced multigraph is connected.  Equip this finite set with uniform
measure.  Every connected configuration has
exactly
\[
  2^{3N/2}(3N/2)!
\]
preimages, obtained by labelling and orienting its edges.  If $q$
sends a point of $\widetilde\Omega_N$ to its unlabelled graph class, then the
pushforward of uniform measure by $q$ is $\pi_N$.  Consequently
\[
 U:L^2(\pi_N)\longrightarrow L^2(\widetilde\Omega_N),
 \qquad (UF)(\omega)=F(q(\omega)),
\]
is an isometric embedding.

Fix a labelled edge $e$.  Define $K_e$ on the cover as follows.  If
the two ends of $e$ occupy different vertex boxes, choose one of the two
remaining positions in each box and exchange the edge ends at those positions;
$K_e$ averages the four possible exchanges.  Each exchange is an involution,
so $K_e$ is self-adjoint for uniform measure.  If $e$ is a loop, set $K_e$ to
be the identity.

For a function in the range of $U$, condition on the positions of the two
ends of $e$ and on all positions outside its endpoint boxes.  On the event that
$e$ is not a loop, these data divide the cover into fibres.  Within one fibre,
the four remaining edge ends can be assigned to the four remaining positions in
$4!$ ways.  Regard these positions as the vertices of $K_{2,2}$, with one
part in each endpoint box.  The allowed exchanges are the transpositions along
its four edges.  Transpositions along the edges of any connected graph
generate the full symmetric group on its vertices, so these four exchanges
connect all $4!$ assignments.  Starting from a connected assignment, every
assignment reached in this way remains connected because each exchange is a
Whitehead move.

Label the four remaining ends $1,2,3,4$.  A function on isomorphism classes is
invariant under permutations of the two positions within either box, so its
value depends only on the two-subset placed in the first box.  The six
possibilities are
\[
  12,\ 13,\ 14,\ 23,\ 24,\ 34.
\]
Each has four preimages among the $4!$ assignments and therefore has uniform
measure.  An exchange replaces one member of the two-subset by one member of
its complement.  Thus two states are adjacent exactly when their intersection
has size one.  The induced operator is simple random walk on the Johnson graph
$J(4,2)$, which is the graph of an octahedron.

Let $A$ be its adjacency matrix.  The three complementary pairs are
$12\leftrightarrow34$, $13\leftrightarrow24$, and
$14\leftrightarrow23$.  Let $C$ be the involution that exchanges each pair.
Then
\[
  A=J-I-C,
\]
where $J$ is the all-ones operator.  The constant vector has eigenvalue $4$.
On its orthogonal complement $J=0$, so the $+1$ and $-1$ eigenspaces of $C$
give eigenvalues $-2$ and $0$, respectively.  The $+1$ eigenspace has
dimension three and contains the constant vector; the $-1$ eigenspace has
dimension three.  Hence $A/4$ has eigenvalues
\[
  1\quad\text{(once)},\qquad
  -\tfrac12\quad\text{(twice)},\qquad
  0\quad\text{(three times)}.
\]
The loop fibres contribute only the eigenvalue $1$.  Averaging the conditional
quadratic forms over all fibres gives, for every $F\in L^2(\pi_N)$,
\[
  -\tfrac12\langle UF,UF\rangle
  \leq\langle UF,K_eUF\rangle
  \leq\langle UF,UF\rangle.
\]

The lifted transition operator is
\[
 \widetilde P_N=\frac{2}{3N}\sum_e K_e.
\]
For a non-loop $e$, two of its four exchanges give each Whitehead resolution;
for a loop, all four are the identity.  Thus the definition of $P_N$ gives the
intertwining identity $\widetilde P_NU=UP_N$.  Averaging the preceding bounds
over $e$ and using that $U$ is an isometry yields
\[
 -\tfrac12\|F\|_{\pi_N}^2
 \leq\langle F,P_NF\rangle_{\pi_N}
 \leq\|F\|_{\pi_N}^2.
\]
Since $P_N$ is self-adjoint, the Rayleigh-quotient characterization of its
extreme eigenvalues proves the proposition.
\end{proof}

We next use the number of loops to bound eigenvalues near $1$.

\section{Loop Count and Low Eigenvalues}\label{sec:loops}

For a reversible kernel, write
\begin{equation}\label{eq:dirichlet}
  \mathcal E_N(F,F)
  =\langle F,(I-P_N)F\rangle_{\pi_N}
  =\frac12\sum_{G,H}\pi_N(G)P_N(G,H)
       \bigl(F(G)-F(H)\bigr)^2.
\end{equation}
Let $L(G)$ denote the number of loops in $G$, and put
$\mu_N(k)=\pi_N\{G:L(G)=k\}$.

\begin{lemma}[Moves that remove a loop]\label{lem:loop-moves}
Let $N\geq4$, and let $G$ have $k$ loops.  Exactly $2k$ Whitehead moves lead
to a graph with $k-1$ loops.  Every move changes the number of loops by at
most one.
\end{lemma}

\begin{proof}
A vertex $u$ carrying a loop $\ell$ has one remaining incident edge $s$, its
\emph{stem}.  Write $v$ for the other endpoint of $s$, and $c,d$ for the two
noncentral ends at $v$.  The two resolutions of $s$ separate the two ends of
$\ell$ and therefore destroy it; see \cref{fig:loop-resolution}.  They destroy
a second loop only if $c$ and $d$ are paired to each other.  In that case both
$u$ and $v$ carry loops and $s$ uses their remaining degrees, so connectedness
forces $N=2$.  Nor can either resolution create a loop: such a loop would have
to come from an edge parallel to $s$, but the two noncentral ends at $u$ are
already paired to form $\ell$.

No other chosen edge can destroy $\ell$: a move on $\ell$ is trivial,
and only its stem can move one of its ends to another vertex.  Distinct loop
vertices have distinct stems when $N\geq4$, since a shared stem would again be
the entire two-vertex component.  Thus the two stem resolutions associated
with each loop are distinct and give exactly $2k$ moves that remove a loop.

Finally, a resolution can create a loop only by putting the two ends of an
edge parallel to the chosen central edge into the same new vertex.  Creating
two loops at once would require both pairs of noncentral ends to be paired
across the two central vertices.  Those two edges together with the
central edge would form an isolated triple edge, again forcing $N=2$.
Applying the same argument to the inverse resolution also rules out the
destruction of two loops.  Hence $|L(G')-L(G)|\leq1$.
\end{proof}

\begin{figure}[ht]
  \centering
  \includegraphics[width=.82\textwidth]{fig-loop-change.pdf}
  \caption{Resolving the stem $s$ of a loop $\ell$.  In either resolution the
  two ends of $\ell$ move to different vertices, while the two other incident
  half-edges separate.  Thus each loop contributes exactly two moves that
  decrease the number of loops.}
  \label{fig:loop-resolution}
\end{figure}

\begin{theorem}[Loop-count Dirichlet form]\label{thm:loop-form}
For $N\geq4$ and every function $f$ on the nonnegative integers,
\begin{equation}\label{eq:loop-form}
  \mathcal E_N(f(L),f(L))
  =\frac{2}{3N}\sum_{k\geq1}k\mu_N(k)
       \bigl(f(k)-f(k-1)\bigr)^2.
\end{equation}
Consequently, if
$0=\lambda_1\leq\lambda_2\leq\cdots$ are the eigenvalues of $I-P_N$, then
for every fixed $r\geq1$,
\begin{equation}\label{eq:eigenvalue-bound}
  \lambda_{r+1}\leq\frac{2r}{3N}+o(N^{-1}).
\end{equation}
\end{theorem}

\begin{proof}
By \cref{lem:loop-moves}, the stationary directed flow from level $k$ to level
$k-1$ is
\[
  \mu_N(k)\frac{2k}{3N}.
\]
Detailed balance says that the reverse aggregate flow has the same value.
The two directions of each transition cancel the factor $1/2$ in
\eqref{eq:dirichlet}.  Summing over $k$ proves \eqref{eq:loop-form}.

It remains to choose a test space.  Let $\nu$ be the $\Pois(1)$ law and define
the Charlier polynomials $C_j$ by the generating function
\cite{KoekoekLeskySwarttouw2010}
\begin{equation}\label{eq:charlier}
  \sum_{j\geq0}C_j(k)\frac{t^j}{j!}=e^{-t}(1+t)^k.
\end{equation}
If $X\sim\nu$, orthogonality follows directly by multiplying two generating
functions:
\[
 \E\!\left[e^{-s}(1+s)^X e^{-t}(1+t)^X\right]
 =e^{-s-t}\exp\bigl((1+s)(1+t)-1\bigr)=e^{st}.
\]
Comparing coefficients, and subtracting the generating functions at $k$ and
$k-1$, gives
\[
  \E[C_i(X)C_j(X)]=j!\,\mathbf 1_{\{i=j\}},
  \qquad
  C_j(k)-C_j(k-1)=jC_{j-1}(k-1).
\]
Since $k\nu(k)=\nu(k-1)$, the limiting form
\[
  \mathcal D(f,f)=\sum_{k\geq1}k\nu(k)
      \bigl(f(k)-f(k-1)\bigr)^2
\]
satisfies
\[
  \mathcal D(C_i,C_j)=j\,j!\,\mathbf 1_{\{i=j\}}.
\]
Fix $r$.  For $0\leq i,j\leq r$, define the Gram and difference matrices
\[
 B_N(i,j)=\sum_{k\geq0}\mu_N(k)C_i(k)C_j(k),
\]
and
\[
 D_N(i,j)=\sum_{k\geq1}k\mu_N(k)
   \bigl(C_i(k)-C_i(k-1)\bigr)
   \bigl(C_j(k)-C_j(k-1)\bigr).
\]
Every entry is the expectation of a fixed polynomial in $L$.  Moment
convergence in \cref{lem:configuration-estimates} therefore gives
\[
 B_N\longrightarrow B=\operatorname{diag}(0!,1!,\ldots,r!),
 \qquad
 D_N\longrightarrow D=\operatorname{diag}(0\,0!,1\,1!,\ldots,r\,r!).
\]
The limiting Gram matrix $B$ is positive definite.  Hence $B_N$ is positive
definite for all sufficiently large $N$, the functions
$C_0(L),\ldots,C_r(L)$ are then linearly independent, and
\[
 \lambda_{\max}\bigl(B_N^{-1/2}D_NB_N^{-1/2}\bigr)
 \longrightarrow
 \lambda_{\max}\bigl(B^{-1/2}DB^{-1/2}\bigr)=r.
\]
Equivalently, the largest quotient $x^{\mathsf T}D_Nx/(x^{\mathsf T}B_Nx)$
on their span is $r+o(1)$.  By \eqref{eq:loop-form}, the largest Laplacian
Rayleigh quotient on that $(r+1)$-dimensional span is therefore
$2r/(3N)+o(N^{-1})$.  The Courant--Fischer min--max principle gives
\eqref{eq:eigenvalue-bound}.
\end{proof}

For $r=1$, one may take the centered test function $L-\E_{\pi_N}L$; its
Rayleigh quotient is $2/(3N)+o(N^{-1})$.  The higher Charlier polynomials
provide the remaining test functions.  The process $L(G_t)$ need not itself
be Markov: transition probabilities between loop counts can depend on the
rest of the graph.  The identity \eqref{eq:loop-form} uses only the total
stationary flow between successive values of $L$.

\section{Bridges and Conductance}\label{sec:bridges}

Let $\mathcal B_N\subset\C_N$ be the classes with at least one bridge, let
$\mathcal B_0=\C_N\setminus\mathcal B_N$, and let $\mathcal B_1$ be the
classes with exactly one bridge.  For sets $S,T\subseteq\C_N$, write
\[
  Q_N(S,T)=\sum_{G\in S,\,H\in T}\pi_N(G)P_N(G,H).
\]
The conductance of the cut between bridged and bridgeless graphs is
\[
  \Phi_N(\mathcal B_N)=
  \frac{Q_N(\mathcal B_N,\mathcal B_0)}
       {\min\{\pi_N(\mathcal B_N),\pi_N(\mathcal B_0)\}}.
\]

We first determine which Whitehead moves cross this cut.  Recall that an edge
is a bridge if and only if it lies on no cycle.

\begin{lemma}[Whitehead moves and bridges]\label{lem:bridge-surgery}
Let $G'$ be obtained from $G$ by resolving an edge $e$.
\par\smallskip
\begin{enumerate}
\item Every bridge of $G$ other than $e$ remains a bridge of $G'$.
\item If $e$ is the unique bridge of $G$, then either resolution of $e$ is
bridgeless.
\end{enumerate}

Consequently, a resolution crosses from $\mathcal B_N$ to $\mathcal B_0$ if
and only if it is one of the two resolutions of the unique bridge of a graph
in $\mathcal B_1$.
\end{lemma}

\begin{proof}
Let $f\ne e$ be a bridge of $G$, and consider the two components of $G-f$.
The edge $e$ lies entirely in one component.  Collapsing and re-expanding
$e$ changes incidences only on that side of the cut: it creates no second
edge between the two components.  Thus $f$ is still their only connection in
$G'$ and remains a bridge.  This proves the first assertion and also shows
that a crossing move must act on the unique bridge.

For the converse, suppose that $e=uv$ is the unique bridge, and let $A$ and
$B$ be the components of $G-e$ containing $u$ and $v$.  Neither side contains
a bridge: an edge separating $A$ or $B$ would also separate $G$.  In one of
the two resolutions, let $a=ux$ be the edge whose end at $u$ moves
to $v$, and let $b=vy$ be the edge whose end at $v$ moves to $u$.
Thus, after the resolution, the incidences of the three crossing edges are
\[
 e=uv,\qquad a=vx,\qquad b=uy.
\]
The other resolution has the same description after choosing the other
incident edge on each side.  This notation also covers a loop: if $a$ was a
loop at $u$, then $x=u$, and moving one of its ends turns it into the edge
$vu$.

Because $a$ is not a bridge of $A$, the graph $A-a$ contains a path $P_A$
from $u$ to $x$; similarly, $B-b$ contains a path $P_B$ from $v$ to $y$.
The three crossing edges lie on the following cycles in $G'$:
\[
 e+a+P_A,\qquad e+b+P_B,\qquad a+P_B+b+P_A.
\]
Here an expression records the edges and paths in their cyclic order; it also
covers the case in which one of $P_A,P_B$ has length zero.

It remains to check edges internal to the two sides.  Let $f$ be an edge of
$A-a$.  Since $A$ is bridgeless, $f$ lies on a cycle $C$ in $A$.  If $C$
avoids $a$, it survives unchanged.  If $C$ uses the old edge $a=ux$, replace
that edge in $C$ by the path $u-e-v-a-x$.  This gives a cycle through $f$ in
$G'$.  The same argument in $B$, replacing the old $b=vy$ by
$v-e-u-b-y$, treats every edge of $B-b$.  Hence every edge of $G'$ lies on a
cycle, and $G'$ is bridgeless.
\end{proof}

\begin{figure}[ht]
  \centering
  \includegraphics[width=.94\textwidth]{fig-bridge-resolution.pdf}
  \caption{One resolution of the unique bridge $e=uv$.  The edge $a=ux$
  becomes $a=vx$, and $b=vy$ becomes $b=uy$.  Since $A-a$ and $B-b$ remain
  connected, their internal paths close cycles through all three crossing
  edges and replace any side cycle that formerly used $a$ or $b$.}
  \label{fig:bridge-resolution}
\end{figure}

\begin{theorem}[Bridge flow]\label{thm:bridge-flow}
For every even $N\geq2$,
\begin{equation}\label{eq:exact-bridge-flow}
  Q_N(\mathcal B_N,\mathcal B_0)
  =\frac{2}{3N}\pi_N(\mathcal B_1).
\end{equation}
Moreover,
\begin{equation}\label{eq:bridge-masses}
  \pi_N(\mathcal B_0)\longrightarrow e^{-1},
  \qquad
  \pi_N(\mathcal B_1)\longrightarrow e^{-1},
  \qquad
  \pi_N(\mathcal B_N)\longrightarrow1-e^{-1},
\end{equation}
and hence
\begin{equation}\label{eq:bridge-conductance}
  \Phi_N(\mathcal B_N)=\frac{2}{3N}(1+o(1)).
\end{equation}
\end{theorem}

\begin{proof}
By \cref{lem:bridge-surgery}, exactly two moves take a graph in
$\mathcal B_1$ to $\mathcal B_0$, whereas no move takes a graph with more
than one bridge to $\mathcal B_0$.  Each move has probability $1/(3N)$.
Summing with respect to $\pi_N$ gives \eqref{eq:exact-bridge-flow}.

In a connected cubic multigraph with $N\geq4$, every loop vertex is joined to
the rest of the graph by a bridge.  Conversely, a bridge that cuts off a single
vertex must have a loop at that vertex.  Each side of a bridge has odd order,
so, if no bridge has at least three vertices on its smaller side, these
singleton-side bridges are all the bridges.  By
\cref{lem:configuration-estimates}, the exceptional event has probability
$O(N^{-1})$ under $\pi_N$.  It follows that $\mathcal B_0$ differs from
$\{L=0\}$ by a set of stationary mass $O(N^{-1})$, and that
$\mathcal B_1$ differs from $\{L=1\}$ by such a set.  The Poisson limit for
$L$ proves the first two limits in \eqref{eq:bridge-masses}; the third follows
because $\mathcal B_N$ is the complement of $\mathcal B_0$.
Since $e^{-1}<1-e^{-1}$, $\mathcal B_0$ is the smaller
side for all sufficiently large $N$.  Substitution in
\eqref{eq:exact-bridge-flow} gives \eqref{eq:bridge-conductance}.
\end{proof}

\begin{theorem}[Exact bridge conductance]\label{thm:exact-conductance}
For every even $N\geq2$, writing $E_N=3N/2$, we have
\begin{equation}\label{eq:bridge-mass-ratio}
 (E_N-1)\pi_N(\mathcal B_1)=E_N\pi_N(\mathcal B_0).
\end{equation}
The bridgeless set $\mathcal B_0$ is the smaller side of the cut, and
\begin{equation}\label{eq:finite-identity}
 \Phi_N(\mathcal B_N)=\frac{1}{E_N-1}.
\end{equation}
\end{theorem}

\begin{proof}
Let $a_G$ be the order of the full incidence automorphism group, including
loop flips and permutations of parallel edges.  Thus
\[
 a_G=|\Aut_v(G)|\prod_i2^{m_{ii}}m_{ii}!\prod_{i<j}m_{ij}!,
\]
and the stationary weight is proportional to $a_G^{-1}$.  Define
$B(t)=\sum_Gt^{|V(G)|}/a_G$ over connected bridgeless cubic graphs, and
$D(t)$ over connected cubic graphs with exactly one bridge.  Let $A(t)$
count connected cubic graphs with one external half-edge and no internal
bridge, with automorphisms fixing that half-edge.  Deleting the unique
bridge gives an unordered pair of graphs counted by $A$.  Removing the
external half-edge and suppressing its vertex gives a bridgeless cubic
graph with a distinguished edge, except for the one-loop graph of weight
$t/2$.  Accounting for symmetry factors gives
\begin{equation}\label{eq:bridge-series-decomposition}
 D=\frac{A^2}{2},\qquad A=\frac t2+\frac32t^2B'.
\end{equation}

All series below are formal.  For a standard Gaussian variable $X$, set
\[
 W(t,J)=\log\E e^{tX^3/6+JX},\qquad
 \phi=W_J(t,J),\qquad \Gamma(t,\phi)=J\phi-W(t,J).
\]
The combinatorial Legendre transform expresses $\Gamma$ in terms of
connected graphs without internal bridges; see
\cite[Section~5.8]{Borinsky2018}.  With the present sign convention its
constant and linear terms are
\[
 \Gamma(t,0)=-B(t),\qquad \Gamma_\phi(t,0)=-A(t).
\]
These terms include vacuum graphs and the one-loop one-point graph.
For completeness, the cubic change of variables gives
\begin{equation}\label{eq:cubic-effective-action}
 \Gamma(t,\phi)=\frac{\phi^2}{2}-\frac{t\phi^3}{6}
 +\frac12\log(1-t\phi)
 -B\bigl(t(1-t\phi)^{-3/2}\bigr).
\end{equation}
Indeed, put $m=1-t\phi$, $s=tm^{-3/2}$, and
$K=(J-\phi+t\phi^2/2)/\sqrt m$.  Formal Gaussian translation and scaling
give
\[
 e^{W(t,J)}=m^{-1/2}
 e^{-\phi^2/2+t\phi^3/6+J\phi}e^{W(s,K)}.
\]
The condition $W_J(t,J)=\phi$ is equivalent to $W_J(s,K)=0$, so
$W(s,K)=B(s)$, proving \eqref{eq:cubic-effective-action}.
This identity is also given in \cite[Section~7.5]{Borinsky2018}.

Twice differentiating \eqref{eq:cubic-effective-action} at $\phi=0$ and
using \eqref{eq:bridge-series-decomposition} yields
\[
 \Gamma_{\phi\phi}(t,0)=1+\frac t2A-\frac{3t^2}{2}A'.
\]
At the zero-mean source $J=-A(t)$, Gaussian integration by parts and the
inverse-Hessian identity give
\[
 A=\frac t2\langle X^2\rangle_{t,-A(t)}
   =\frac{t}{2\Gamma_{\phi\phi}(t,0)}.
\]
Here the brackets denote the normalized formal expectation.  Consequently,
\[
 3tAA'-A^2=\frac{2A}{t}-1,
 \qquad
 \left(\frac32t\frac{d}{dt}-1\right)D
 =\frac32tB'.
\]
Taking the coefficient of $t^N$ proves
$(E_N-1)[t^N]D=E_N[t^N]B$, hence \eqref{eq:bridge-mass-ratio}.
Bridgeless cubic graphs exist at every even order, for instance a cycle
with a perfect matching added.  Since $\mathcal B_1\subseteq\mathcal B_N$
and $E_N/(E_N-1)>1$, the bridgeless side is strictly smaller.  The exact flow
formula \eqref{eq:exact-bridge-flow} now gives \eqref{eq:finite-identity}.
\end{proof}

We use the usual worst-case total-variation mixing time
\[
  t_{\mathrm{mix}}(\varepsilon)=
  \min\bigl\{t:\max_G\|P_N^t(G,\cdot)-\pi_N\|_{\TV}\leq\varepsilon\bigr\}.
\]
We define the relaxation time by
$t_{\rm rel}=1/\lambda_2(I-P_N)$.

\begin{corollary}[Mixing lower bound]\label{cor:mixing}
For every even $N\geq2$ and $0<\varepsilon<1$,
\begin{equation}\label{eq:finite-mixing-lower}
 t_{\mathrm{mix}}(\varepsilon)
 \geq(3N/2-1)\bigl(1-\pi_N(\mathcal B_0)-\varepsilon\bigr).
\end{equation}
In particular, for every fixed $\varepsilon<1-e^{-1}$,
\begin{equation}\label{eq:mixing-lower}
  t_{\mathrm{mix}}(\varepsilon)
  \geq \frac{3N}{2}
      \bigl(1-e^{-1}-\varepsilon\bigr)(1+o(1)).
\end{equation}
In particular,
\[
  t_{\mathrm{mix}}(1/4)\geq0.57318\ldots\,N(1+o(1)).
\]
The relaxation time is also at least $(3/2+o(1))N$.
\end{corollary}

\begin{proof}
Start the chain from $\mu_0=\pi_N(\,\cdot\mid\mathcal B_0)$.  Pointwise,
$\mu_0\leq\pi_N/\pi_N(\mathcal B_0)$, and applying a Markov kernel preserves
this domination.  If $\mu_s=\mu_0P_N^s$, the probability of crossing from
$\mathcal B_0$ to $\mathcal B_N$ at step $s+1$ is therefore
\[
 \sum_{G\in\mathcal B_0}\mu_s(G)P_N(G,\mathcal B_N)
 \leq
 \frac{Q_N(\mathcal B_0,\mathcal B_N)}{\pi_N(\mathcal B_0)}.
\]
A union bound over the first $t$ steps gives
\[
  \mu_0P_N^t(\mathcal B_0)
  \geq1-t\frac{Q_N(\mathcal B_0,\mathcal B_N)}
                       {\pi_N(\mathcal B_0)}
  =1-t\Phi_N(\mathcal B_N)
\]
for every even $N\geq2$; the equality uses reversibility and
\cref{thm:exact-conductance}.  Testing total variation on
$\mathcal B_0$ gives
\[
 \|\mu_0P_N^t-\pi_N\|_{\TV}
 \geq1-t\Phi_N(\mathcal B_N)-\pi_N(\mathcal B_0).
\]
Using \eqref{eq:finite-identity} proves \eqref{eq:finite-mixing-lower}, and
the limit in \eqref{eq:bridge-masses} gives \eqref{eq:mixing-lower}.
Finally, convexity of total variation gives
\[
 \|\mu_0P_N^t-\pi_N\|_{\TV}
 \leq\sum_G\mu_0(G)\|P_N^t(G,\cdot)-\pi_N\|_{\TV}
 \leq\max_G\|P_N^t(G,\cdot)-\pi_N\|_{\TV},
\]
so the lower bound obtained from the conditional initial law also holds for
worst-case mixing.
The relaxation claim follows from \eqref{eq:eigenvalue-bound} with $r=1$:
$\lambda_2(I-P_N)\leq2/(3N)+o(N^{-1})$, and the reciprocal is the relaxation
time.  See \cite{LevinPeres2017} for the standard mixing and conductance
definitions.
\end{proof}

\section{Enumeration}\label{sec:enumeration}

For fixed $N$, we enumerate the isomorphism classes in $\C_N$ and compute
the integer matrix $M_N=(m_{GH})$ by a breadth-first search.  Every Whitehead
move is counted, including moves with isomorphic results and those that leave
the class unchanged.  The tables below are computed from these matrices.

\subsection{Canonical Representation}

Store a cubic multigraph by its upper-triangular adjacency-multiplicity
matrix, with diagonal entries counting loops.  To compare two such matrices,
form a colored simple incidence graph.  Its first color class consists of
the $N$ original vertices.  Every non-loop edge becomes a vertex of a
second color, adjacent to its two endpoints.  Every loop becomes a
vertex of a third color, adjacent to its single incident original vertex.
Keeping loops and non-loop edges in different colors retains the
information lost by replacing a loop's double incidence with one simple
incidence.
Parallel edges remain distinct incidence vertices.

Two multigraphs are isomorphic if and only if their colored incidence graphs
are color-preservingly isomorphic.  Canonical representation is a standard
way to avoid duplicate isomorphism classes in exhaustive generation
\cite{McKay1998}.  We use nauty/Traces \cite{McKayPiperno2014} to canonically
label the incidence graph.  Reading the original-vertex color class in
canonical order and writing the corresponding upper triangle produces a
canonical code.  Two codes are equal exactly when the corresponding
multigraphs are isomorphic.

\subsection{Whitehead Moves and Breadth-First Search}

The program represents each of the $3N/2$ edges by a pair of half-edges
and records the vertex incident with each half-edge.  For a non-loop edge
$e=uv$, let $s_u,s_v$ be its half-edges, let $a,b$ be the other half-edges
at $u$, and let $c,d$ be the other half-edges at $v$.  The pairing is
unchanged by a Whitehead move.  The two resulting graphs are obtained by
changing the vertex incidences as follows:
\begin{align*}
 &a,c\text{ are incident with }u,
 &&b,d\text{ are incident with }v;\\
 &a,d\text{ are incident with }u,
 &&b,c\text{ are incident with }v.
\end{align*}
The half-edges $s_u,s_v$ remain at $u,v$.  Converting the resulting incidence map
back to an adjacency-multiplicity matrix implements exactly the two
partitions in \cref{fig:whitehead-move}, including cases in which two of
$a,b,c,d$ belong to the same edge.

\begin{algorithm}[Transition enumeration]\label{alg:enumeration}
Fix an even $N$ and a connected cubic starting graph.  Maintain a set
$\mathcal S$ of canonical codes and a first-in, first-out queue, initially
containing only the starting code.  While the queue is nonempty:

\begin{enumerate}
\item remove one code and reconstruct its adjacency multigraph;
\item expand its $3N/2$ edges into paired half-edges, keeping parallel
edges distinct;
\item for a loop, increase the diagonal entry of $M_N$ by two; for a non-loop
edge, perform the two Whitehead moves described above;
\item canonicalize each non-loop target, increment its transition
multiplicity, and enqueue its code if it is not already in $\mathcal S$.
\end{enumerate}

Return $\mathcal S$ and the integer multiplicity matrix
$M_N=(m_{GH})$.
\end{algorithm}

The search starts from the triple edge when $N=2$.  For $N\geq4$, take a path
$v_1v_2\cdots v_N$, put a loop at each endpoint, and, for $1\leq i<N$, give
$v_iv_{i+1}$ multiplicity one when $i$ is odd and two when $i$ is even.
Adjacent path edges then have total multiplicity three at every internal
vertex, while each endpoint has its loop and one path edge.

\begin{proposition}[Correctness of the enumeration]
\Cref{alg:enumeration} terminates and returns every class in $\C_N$ exactly
once.  For every ordered pair of returned classes, its matrix entry is the
number of Whitehead moves between them.
\end{proposition}

\begin{proof}
There are finitely many adjacency-multiplicity matrices with $N$ vertices and
degree three, so the queue terminates.  Canonical codes identify precisely
the isomorphism classes, which proves that no class is stored twice.  Every
generated graph is cubic because $u$ and $v$ each remain incident with three
half-edges.  It is connected because collapsing $e$ gives a connected graph and
either displayed expansion reconnects the two new vertices by $e$.  The
Whitehead support on $\C_N$ is connected by
\cref{prop:connected-support}; hence a shortest support path from the starting
class to any class is discovered inductively by the breadth-first search.
Finally, the inner loop visits each edge and both Whitehead moves on it.
Moves with the same resulting class are counted separately, giving the
multiplicities in \eqref{eq:kernel}.
\end{proof}

Let $C_N=|\C_N|$, and let $T_{\rm can}(N)$ denote the maximum time for one
canonical labelling of the colored incidence graphs used here.  Such a graph
has $5N/2$ vertices and at most $3N$ incidence edges.  The search generates
exactly $3NC_N$ Whitehead moves and makes at most $3NC_N+1$
canonical-labelling calls, including the starting code.  For each of the $C_N$
states, constructing and copying the dense $N\times N$ target matrices for
$3N$ moves costs $O(N^3)$ time.  Canonical
labelling contributes at most $3N T_{\rm can}(N)$ per state.  The total time is
\[
 O\!\left(C_N\bigl(N^3+N T_{\rm can}(N)\bigr)\right).
\]
Storing a dense code for each state costs $O(N^2C_N)$ machine words.  If
transitions are stored as pairs of state indices with multiplicities, they
use $O(NC_N)$ more, giving $O(N^2C_N)$ total storage.  The supplied program
instead retains a separate record for every Whitehead move, including its
target code and vertex permutations.  Its storage bound is therefore
$O(N^3C_N)$.

\subsection{Computational Results}

All counts are computed over the integers.  For each $N$, the number of
classes agrees with OEIS A005967 \cite{OEISA005967}, every row of $M_N$ sums
to $3N$, and every transition between distinct classes has a reverse.
Stationary weights are obtained by propagating the rational identities
$w_Gm_{GH}=w_Hm_{HG}$ along the connected support and checking them on every
edge.  Normalizing these weights gives the stationary masses, cut flow, and
conductance as reduced fractions.  This calculation uses only $M_N$, not the
formula in \cref{thm:stationary}; it checks the table's stationary quantities
without separately verifying the automorphism factors in that formula.

The bridges of each graph are found by deleting each edge and testing
connectedness.  These counts confirm \cref{lem:bridge-surgery}: a class with
one bridge has exactly two moves into $\mathcal B_0$, and a class with more
than one bridge has none.  The conductance is computed from the transition
matrix before comparison with \eqref{eq:finite-identity}.

The computations used the native Traces backend with one worker on an Apple
M5 computer with 24 GB of memory.  At $N=14$, the search visited $21{,}096$
classes and all $886{,}032$ Whitehead moves.  The accompanying record contains
the counts and rational values, running times, software versions, source
hashes, and SHA-256 digests of the ordered canonical codes and sparse
transition matrix.  The full matrices can be reconstructed by the supplied
programs.  The tables are generated from this record.  In
\cref{tab:enumeration-size}, support edges are unordered pairs of distinct
classes joined by a Whitehead move.

% Generated by render_rule_a_table.py from rule-a-enumeration.json.
% Do not edit this file by hand.
\newcommand{\RuleASizeRows}{%
2 & 2 & 1 & 12 & $<0.01$ \\
4 & 5 & 4 & 60 & $<0.01$ \\
6 & 17 & 25 & 306 & 0.01 \\
8 & 71 & 179 & 1{,}704 & 0.14 \\
10 & 388 & 1{,}537 & 11{,}640 & 1.65 \\
12 & 2{,}592 & 15{,}402 & 93{,}312 & 17.67 \\
14 & 21{,}096 & 177{,}630 & 886{,}032 & 238.6 \\
}
\newcommand{\RuleABridgeRows}{%
2 & $\dfrac{2}{5}$ & $\dfrac{3}{5}$ & $\dfrac{1}{5}$ & $\dfrac{1}{2}$ \\
4 & $\dfrac{1}{3}$ & $\dfrac{2}{5}$ & $\dfrac{1}{15}$ & $\dfrac{1}{5}$ \\
6 & $\dfrac{352}{1\,105}$ & $\dfrac{396}{1\,105}$ & $\dfrac{44}{1\,105}$ & $\dfrac{1}{8}$ \\
8 & $\dfrac{539}{1\,695}$ & $\dfrac{196}{565}$ & $\dfrac{49}{1\,695}$ & $\dfrac{1}{11}$ \\
10 & $\dfrac{133\,392}{414\,125}$ & $\dfrac{28\,584}{82\,825}$ & $\dfrac{9\,528}{414\,125}$ & $\dfrac{1}{14}$ \\
12 & $\dfrac{6\,443}{19\,675}$ & $\dfrac{6\,822}{19\,675}$ & $\dfrac{379}{19\,675}$ & $\dfrac{1}{17}$ \\
14 & $\dfrac{85\,315\,584}{256\,406\,305}$ & $\dfrac{447\,906\,816}{1\,282\,031\,525}$ & $\dfrac{21\,328\,896}{1\,282\,031\,525}$ & $\dfrac{1}{20}$ \\
}


\begin{table}[ht]
\centering
\caption{Enumeration sizes and running times.  Times include
generation and verification with one worker on the machine described in the
text; table rendering is excluded.  The first row also includes loading the
native library.}
\label{tab:enumeration-size}
\begin{tabular}{@{}r r r r r@{}}
\toprule
$N$ & $|\C_N|$ & support edges & moves & seconds \\
\midrule
\RuleASizeRows
\bottomrule
\end{tabular}
\end{table}

\begin{table}[ht]
\centering
\small
\setlength{\tabcolsep}{7pt}
\renewcommand{\arraystretch}{2.2}
\caption{Stationary bridge data computed from the full transition matrices.
Here $\mathcal B_0$ is the bridgeless set, $\mathcal B_1$ is the
one-bridge set, and $\mathcal B_N$ is the set of all bridged classes.}
\label{tab:enumeration-bridge}
\begin{tabular}{@{}r c c c c@{}}
\toprule
$N$ & $\pi_N(\mathcal B_0)$ & $\pi_N(\mathcal B_1)$
& $Q_N(\mathcal B_N,\mathcal B_0)$ & $\Phi_N(\mathcal B_N)$ \\
\midrule
\RuleABridgeRows
\bottomrule
\end{tabular}
\end{table}

The last column of \cref{tab:enumeration-bridge} verifies
\cref{thm:exact-conductance} for all seven enumerated sizes.  The
corresponding configuration-count identity is obtained by setting
\[
 W_i=\sum_{G\in\mathcal B_i}
       \frac{N!}{|\Aut_v(G)|}c(G),
 \qquad i=0,1.
\]
By \cref{lem:configuration-fibers}, $W_i$ is the number of connected
half-edge configurations whose unlabelled graph has $i$ bridges.  Then
\begin{equation}\label{eq:finite-weight-identity}
  (3N/2-1)W_1=(3N/2)W_0.
\end{equation}

This is \eqref{eq:bridge-mass-ratio} multiplied by the common stationary
normalizing factor, and holds for every even $N\geq2$.

The programs and data used to produce
\cref{tab:enumeration-size,tab:enumeration-bridge} accompany the arXiv version
as ancillary files.  Instructions are included for reproducing the
computations for every even $N$ from $2$ through $14$.

\section{Concluding Remarks}

The local Whitehead chain is the random walk on the directed graph of graphs
in which every move is counted with multiplicity.  Its stationary measure is
given by \eqref{eq:stationary}.  The spectral and conductance estimates show
that both relaxation and mixing require at least a linear number of steps.

The bridge cut has the exact conductance \eqref{eq:finite-identity}.
It remains to determine whether there is an absolute
$c>0$ such that
\[
  \lambda_2(I-P_N)\geq \frac{c}{N}.
\]
Together with \cref{thm:loop-form}, such a bound would determine the
order of the relaxation time.  Finally, one may ask whether the Charlier
upper bounds are asymptotically sharp for the individual eigenvalues.

\section*{Declaration of interests}

The author has no competing interests to declare.

\bibliographystyle{abbrv}
\bibliography{local-whitehead-chain}

\end{document}
